\section{Conclusions}
This review chapter delivers a holistic overview of literature pertaining to the spatio-temporal climate sensitivity of the atmosphere to reactive non carbon-based species. This includes literature from a range of research areas: aircraft emissions, plume dispersion and dynamics, the distribution of air traffic and corresponding emissions, the climate impact of aviation on both a global and a local scale, and lastly the mitigation potential for both \ce{CO_2} and non-\ce{CO_2} climate impact. 

A strong finding from this literature review, giving direction to fruitful future work on this research topic, is the effect of plume-scale nonlinear chemistry and microphysics on the eventual climate impact of reactive non-\ce{CO_2} aircraft emissions. The chemical and physical state of the atmosphere is affected by both natural variability in climate parameters, as well as by perturbations to atmospheric chemistry due to anthropogenic emissions. In the case of aircraft operating in the UTLS, the atmospheric state (and hence the chemical fate of reactive aircraft emissions) is influenced by both the background chemical composition of the surrounding air and the elevated concentrations of emissions within the aircraft exhaust plume. The sensitivity of the atmospheric response to emissions is heightened at these altitudes compared to at surface level, due to reasons elaborated on in section \ref{UTLS_climate}. Climate-optimal aircraft routing and formation flight can both be utilised to minimise time spent flying through atmospheric regions that are unfavourable in terms of climate impact.

Real world implementation of non-carbon based operational mitigation concepts does however require: knowledge of aviation fuel consumption, plume dispersion characteristics, and monitoring and forecasting of chemical and meteorological parameters along potential routes. Moreover, this information must be obtained to a sufficiently high-resolution and within a relatively short timespan, so as to capture adequate variability in climate parameters to make an informed and timely decision on routing. As a result, modelling assumptions will inevitably need to be made, to balance computational run time with accuracy. In particular, the incorporation of plume-scale effects into large-scale climate models is likely to bring about issues with resolving data between the two resolutions. Another key issue limiting feasibility at the current time is that of air traffic management. Further work should focus on these key challenges and work towards improving scientific understanding of the non-\ce{CO_2} radiative forcing multiplier. Achieving a general scientific consensus on this non-\ce{CO_2} multiplier would increase confidence amongst policymakers to include such effects in aviation climate policy, and hence accelerate the widespread adoption of interim mitigation approaches such as climate-optimal aircraft routing and formation flight.

The extensive background reading done to write this review chapter is used throughout the thesis, to inform the later chapters and to ensure work done is truly novel and relevant. The rest of this thesis focuses on a number of different studies conducted by my team and I at the University of Bristol, across both Aerospace Engineering and the Atmospheric Chemistry Research Group (ACRG). It is worth noting that this chapter is based on a publication from 2022, meaning it does not address some key papers on aircraft performance, atmospheric chemistry and climate modelling published since then. It is left to the remaining chapters of this thesis to explain the necessary updates to the state-of-the-art science and modelling methods where relevant.